\documentclass[../cours_tle.tex]{subfiles}
\begin{document}
\setcounter{chapter}{0}
\chapter{Langage, fonctions et calcul : les bases}
\label{ch:1}
\begin{objectif}
Spécialité. Revoir ensembles, équations et fonctions avant toute étude de limite. Prérequis : calcul littéral et fonctions de seconde, chapitres 1 à 3 de première.
\end{objectif}
\section{Définir les objets avant de calculer}
\begin{definition}
Un ensemble rassemble des objets appelés éléments. $x\in E$ exprime l'appartenance ; $A\subset E$ signifie que tout élément de $A$ appartient à $E$.
Une fonction $f:E\to F$ associe à chaque $x\in E$ une unique image $f(x)\in F$.
Un antécédent de $y$ est un $x$ tel que $f(x)=y$ : il peut n'y en avoir aucun, un ou plusieurs.
\end{definition}
$\N=\{0,1,2,\ldots\}$, $\Z$ contient aussi les entiers négatifs, $\Q$ les quotients d'entiers de dénominateur non nul ; $\R$ contient notamment $\sqrt2$ et $\pi$.
Un intervalle contient tous les réels situés entre deux de ses éléments. $[a,b]$ inclut ses bornes ; $]a,b[$ les exclut. L'infini n'est pas un réel et ne se met jamais entre crochets fermés.

$A\cap B$ exige les deux appartenances, $A\cup B$ au moins une. Le complémentaire est pris dans un ensemble de référence fixé. Le produit cartésien $A\times B$ est l'ensemble des couples ordonnés $(a,b)$ : $(1,2)\ne(2,1)$.
\section{Logique et négation}
$P\Rightarrow Q$ affirme que $P$ suffit pour $Q$ ; $Q$ est nécessaire à $P$.
La réciproque $Q\Rightarrow P$ n'est pas automatique. La contraposée $\neg Q\Rightarrow\neg P$ est équivalente à l'implication initiale : toutes deux excluent le cas $P$ vrai et $Q$ faux.
$P\Leftrightarrow Q$ exige deux implications.

\begin{exemple}
$x=2\Rightarrow x^2=4$ mais $x^2=4$ permet aussi $x=-2$.
Pour prouver que $n^2$ impair entraîne $n$ impair, on prouve la contraposée : si $n=2k$, alors $n^2=4k^2$ est pair.
\end{exemple}
Les symboles $\forall$ (« pour tout ») et $\exists$ (« il existe »), non exigibles au lycée, rendent les phrases plus courtes. Leur ordre importe : $\forall x\,\exists y:y>x$ est vrai avec $y=x+1$ ; $\exists y\,\forall x:y>x$ est faux.
Nier $\forall x\in E,P(x)$ donne $\exists x\in E,\neg P(x)$.
Nier « il existe $x$ tel que $P$ et $Q$ » donne « pour tout $x$, non $P$ ou non $Q$ ».
Un seul contre-exemple réfute une affirmation universelle ; de nombreux exemples ne la démontrent pas.
\section{Domaine, équivalences et pièges de calcul}
Avant une équation, imposer dénominateurs non nuls, radicandes positifs ou nuls et arguments des logarithmes strictement positifs.
Multiplier une inégalité par un nombre négatif renverse son sens. Élever au carré est réversible pour deux membres positifs ou nuls, mais pas sans condition de signe.
\begin{exemple}
Résoudre $\sqrt{x+2}=x$ impose $x\ge0$. Le carré donne $x^2-x-2=0$, donc $x=2$ ou $-1$ ; seul $2$ convient.
Pour $(x-1)/(x+2)\ge0$, les points critiques sont $-2$ (interdit) et $1$ (autorisé). Le tableau des signes donne $]-\infty,-2[\cup[1,+\infty[$.
\end{exemple}
\begin{propriete}
Pour tout réel $x$, $\sqrt{x^2}=|x|$. Pour $a,b\in\R$, $|a+b|\le|a|+|b|$.
\end{propriete}
\begin{demonstration}
La racine carrée est la solution positive ou nulle de $t^2=x^2$, donc $|x|$.
Comme $-|a|\le a\le|a|$ et de même pour $b$, leur somme est comprise entre $-(|a|+|b|)$ et $|a|+|b|$.
\end{demonstration}
\section{Composition, images et réciproques}
$(g\circ f)(x)=g(f(x))$ n'a de sens que si $f(x)$ appartient au domaine de $g$.
Une fonction est injective si deux images égales imposent deux antécédents égaux ; surjective si chaque élément de l'ensemble d'arrivée possède un antécédent ; bijective si les deux conditions sont réunies.
Alors $f^{-1}$ associe à $y$ son unique antécédent. Ne pas confondre $f^{-1}$ et $1/f$.
Ainsi $x\mapsto x^2$ est bijective de $[0,+\infty[$ sur lui-même, mais pas de $\R$ sur $\R$.

\begin{methode}
Pour une preuve, écrire les données, le résultat visé, puis choisir : calcul direct, distinction de cas, contraposée ou contradiction. Pour une égalité d'ensembles, prouver les deux inclusions en suivant un élément arbitraire.
\end{methode}
\begin{finchapitre}
\section{Utilité et suite du parcours}
La même définition de fonction servira aux suites, aux variables aléatoires et aux applications linéaires. Les domaines et quantificateurs conditionneront toutes les preuves de licence.
\begin{retenir}
Un calcul correct ne suffit pas : préciser les objets, les domaines et le sens de chaque implication.
\end{retenir}
\end{finchapitre}
\end{document}
